\documentclass[12pt]{article}
\usepackage[utf8]{inputenc}
\usepackage[greek,english]{babel}
\usepackage{alphabeta}
\usepackage{graphicx}
\usepackage{amsmath}
\usepackage{amssymb}
\usepackage[hidelinks]{hyperref}
\usepackage{booktabs}
\usepackage{float}
\usepackage{parskip} % Για καλύτερο spacing παραγράφων
\usepackage{listings} % Για κώδικα
\usepackage{xcolor} % Για χρώματα στον κώδικα
\usepackage[top=2cm, bottom=2cm, left=2.5cm, right=2.5cm]{geometry} % Μικρότερα περιθώρια
\newcommand{\en}[1]{\foreignlanguage{english}{#1}}

% Ρυθμίσεις για τον κώδικα Python
\lstset{
    language=Python,
    basicstyle=\ttfamily\small,
    keywordstyle=\color{blue},
    stringstyle=\color{red},
    commentstyle=\color{green!60!black},
    numbers=left,
    numberstyle=\tiny\color{gray},
    breaklines=true,
    frame=single,
    showstringspaces=false
}

\begin{document}
\selectlanguage{greek}

% Ξεχωριστή Σελίδα Τίτλου
\begin{titlepage}
\begin{center}
    % Banner Placeholder
    \vspace*{-1cm}
    \begin{figure}[H]
    \centering
    \foreignlanguage{english}{
\includegraphics[width=0.8\textwidth]{upatras_banner.jpg}}
    \end{figure}
    
    \vspace{0.5cm}
    
    {\large Πανεπιστημιούπολη Ρίου $\cdot$ Πολυτεχνική Σχολή} \\
    {\large Τμήμα Μηχανικών Η/Υ \& Πληροφορικής (\foreignlanguage{english}{ΤΜΗΥΠ} -- \foreignlanguage{english}{CEID})} \\
    
    \vspace{2cm}
    
    {\Huge \textbf{ΤΟΠΟΛΟΓΙΚΗ ΑΝΑΛΥΣΗ}} \\[0.4cm]
    {\Huge \textbf{ΔΕΔΟΜΕΝΩΝ}} \\
    \vspace{0.8cm}
    {\Large (\foreignlanguage{english}{CEID\_25EE606})} \\
    \vspace{0.5cm}
    {\Large Εξαμηνιαία Εργασία --- 2026}
    
    \vspace{2.5cm}
    
    {\large \textbf{Ομάδα Εργασίας:}} \\
    \vspace{0.5cm}
    {\large Βερύκιος Άγγελος (1100500)} \\[0.3cm]
    {\large Βογινατζής Αναστάσιος (1100506)} \\[0.3cm]
    {\large Κολύβρας Κωνσταντίνος (1103826)}
    
    \vfill
    
    {\large 11 Μαΐου 2026}
\end{center}
\end{titlepage}

\newpage
\tableofcontents
\newpage

\section{\texorpdfstring{\underline{Θεωρητική Εισαγωγή}}{Θεωρητική Εισαγωγή}}
\vspace{1em}

Στην ενότητα αυτή παρουσιάζονται οι βασικές έννοιες της Τοπολογικής Ανάλυσης Δεδομένων (\foreignlanguage{english}{TDA}).

\subsection{\foreignlanguage{english}{Simplicial Complexes}}
\vspace{1em}

Τα \foreignlanguage{english}{Simplicial Complexes} αποτελούν τη διακριτή προσέγγιση των συνεχών τοπολογικών χώρων. Αποτελούνται από απλόπλοκα (\foreignlanguage{english}{simplices}) διαφόρων διαστάσεων: 0-απλόπλοκα (σημεία), 1-απλόπλοκα (ακμές), 2-απλόπλοκα (τρίγωνα), 3-απλόπλοκα (τετράεδρα) κ.ο.κ. Η ένωση αυτών των στοιχείων με συγκεκριμένους κανόνες συνδυαστικότητας μας επιτρέπει να μοντελοποιήσουμε το σχήμα των δεδομένων και να υπολογίσουμε τοπολογικές ιδιότητες.

\subsection{\foreignlanguage{english}{Persistent Homology}}
\vspace{1em}

Η \foreignlanguage{english}{Persistent Homology} είναι ένα εργαλείο που ανιχνεύει τοπολογικά χαρακτηριστικά (όπως συνεκτικές συνιστώσες, τρύπες, κοιλότητες) σε πολλαπλές κλίμακες. Καθώς αυξάνουμε την ακτίνα $r$ γύρω από κάθε σημείο (διήθηση ή \foreignlanguage{english}{filtration}), νέα απλόπλοκα δημιουργούνται. Η εμμονή (\foreignlanguage{english}{persistence}) μετρά πόσο "ζει" κάθε τοπολογικό χαρακτηριστικό. Χαρακτηριστικά που επιμένουν για μεγάλο εύρος τιμών του $r$ θεωρούνται αληθή τοπολογικά σήματα, ενώ αυτά που πεθαίνουν γρήγορα θεωρούνται θόρυβος.

\subsection{\foreignlanguage{english}{Mapper}}
\vspace{1em}

Ο αλγόριθμος \foreignlanguage{english}{Mapper} είναι μια τεχνική μείωσης διαστάσεων και οπτικοποίησης που διατηρεί την τοπολογική δομή των δεδομένων. Χρησιμοποιεί μια συνάρτηση φίλτρου (\foreignlanguage{english}{filter function}) για να προβάλλει τα δεδομένα σε χαμηλότερη διάσταση, χωρίζει τον χώρο σε επικαλυπτόμενα διαστήματα, εφαρμόζει \foreignlanguage{english}{clustering} στα σημεία κάθε διαστήματος και συνδέει τα \foreignlanguage{english}{clusters} που μοιράζονται κοινά σημεία, δημιουργώντας ένα γράφημα που συνοψίζει τη δομή.

\newpage
\section{\texorpdfstring{\underline{Μεθοδολογία}}{Μεθοδολογία}}
\vspace{1em}

Για την υλοποίηση της εργασίας ακολουθήθηκε η εξής διαδικασία:

\subsection{Βιβλιοθήκες}
\vspace{1em}

\begin{itemize}
    \item \foreignlanguage{english}{\texttt{GUDHI}}: Για την κατασκευή συμπλεγμάτων και τον υπολογισμό \foreignlanguage{english}{Persistent Homology} στο Ερώτημα Α.
    \item \foreignlanguage{english}{\texttt{Ripser}}: Για τον γρήγορο υπολογισμό της \foreignlanguage{english}{Persistent Homology} στο Ερώτημα Β.
    \item \foreignlanguage{english}{\texttt{Persim}}: Για τον υπολογισμό αποστάσεων \foreignlanguage{english}{Bottleneck} και \foreignlanguage{english}{Wasserstein}.
    \item \foreignlanguage{english}{\texttt{NumPy}} \& \foreignlanguage{english}{\texttt{SciPy}}: Για τη διαχείριση δεδομένων και μαθηματικούς υπολογισμούς.
    \item \foreignlanguage{english}{\texttt{Matplotlib}} \& \foreignlanguage{english}{\texttt{Plotly}}: Για την οπτικοποίηση των αποτελεσμάτων.
\end{itemize}

\subsection{Αιτιολόγηση Επιλογών}
\vspace{1em}

Η επιλογή της \foreignlanguage{english}{GUDHI} έγινε λόγω της πληρότητάς της σε διάφορα είδη συμπλεγμάτων (\foreignlanguage{english}{Alpha Complex}, \foreignlanguage{english}{Cech Complex} κ.λπ.), ενώ η \foreignlanguage{english}{Ripser} επιλέχθηκε για την ταχύτητά της στον υπολογισμό του $H_1$. Ο αλγόριθμος \foreignlanguage{english}{Mapper} υλοποιήθηκε τόσο εκ του μηδενός (για βαθύτερη κατανόηση) όσο και μέσω της βιβλιοθήκης \foreignlanguage{english}{\texttt{kmapper}} για επαλήθευση των αποτελεσμάτων.

\newpage
\section{\texorpdfstring{\underline{Αποτελέσματα \& Ανάλυση}}{Αποτελέσματα \& Ανάλυση}}
\vspace{1em}

\subsection{Ερώτημα Α}
\vspace{1em}

\subsubsection{Α.1—Νέφος Σημείων \& \foreignlanguage{english}{Vietoris--Rips} (\foreignlanguage{english}{from scratch})}
\vspace{0.5em}

\vspace{1em}
\textbf{Εξήγηση Υλοποίησης:}
Για την υλοποίηση εκ του μηδενός, δημιουργήσαμε τα νέφη σημείων και υπολογίσαμε τον πίνακα αποστάσεων. Στη συνέχεια, για μια δεδομένη ακτίνα $r$, βρήκαμε όλες τις ακμές (αποστάσεις $\le 2r$) και όλα τα τρίγωνα (κλίκες μεγέθους 3). Ο υπολογισμός των αριθμών \foreignlanguage{english}{Betti} έγινε μέσω του πίνακα συνόρων (\foreignlanguage{english}{boundary matrix}) και απαλοιφής \foreignlanguage{english}{Gauss}.

\vspace{1em}
\textbf{Σημαντικό Κομμάτι Κώδικα:}
\selectlanguage{english}
\begin{lstlisting}
# Finding edges and triangles (Simplices)
edges = []
for i in range(n):
    for j in range(i+1, n):
        if dist_matrix[i, j] <= 2 * r:
            edges.append((i, j))

triangles = []
for i in range(n):
    for j in range(i+1, n):
        for k in range(j+1, n):
            if dist_matrix[i, j] <= 2*r and dist_matrix[i, k] <= 2*r and dist_matrix[j, k] <= 2*r:
                triangles.append((i, j, k))
\end{lstlisting}
\selectlanguage{greek}
\textit{Σχολιασμός: Ο κώδικας αυτός ελέγχει όλους τους συνδυασμούς σημείων για να βρει αν σχηματίζουν ακμές ή τρίγωνα με βάση την ακτίνα $r$. Είναι η βάση για την κατασκευή του συμπλέγματος.}

\vspace{1em}
\textbf{Αποτελέσματα \& Ερμηνεία:}
Τα αποτελέσματα του υπολογισμού των αριθμών \foreignlanguage{english}{Betti} για τον κύκλο με $\epsilon = 0.5$ έδωσαν $\beta_0 = 1$, $\beta_1 = 1$, $\beta_2 = 1$. \textit{Σημείωση: Η τιμή $\beta_2=1$ οφείλεται στην αποκοπή του συμπλέγματος στη διάσταση 2 (έλλειψη τετραέδρων).}

\begin{figure}[H]
\centering
\foreignlanguage{english}{\includegraphics[width=0.8\textwidth]{point_clouds.png}}
\caption{Τα 4 παραγόμενα νέφη σημείων. Παρατηρούμε τις διαφορετικές τοπολογικές δομές (κύκλος, τόρος, σφαίρα).}
\end{figure}

\subsubsection{Α.2—Σύγκριση \foreignlanguage{english}{VR/Cech/Alpha} μέσω \foreignlanguage{english}{GUDHI}}
\vspace{0.5em}

\vspace{1em}
\textbf{Αποτελέσματα \& Ερμηνεία:}
\begin{table}[H]
\centering
\caption{Σύγκριση Συμπλεγμάτων για $r = 0.3$}
\begin{tabular}{lcccccc}
\toprule
Σύμπλεγμα & \foreignlanguage{english}{0-simplices} & \foreignlanguage{english}{1-simplices} & \foreignlanguage{english}{2-simplices} & $\beta_0$ & $\beta_1$ & $\beta_2$ \\
\midrule
\foreignlanguage{english}{VR Complex} & 30 & 72 & 54 & 1 & 1 & 0 \\
\foreignlanguage{english}{Cech Complex} & 30 & 56 & 26 & 1 & 1 & 0 \\
\foreignlanguage{english}{Alpha Complex} & 30 & 37 & 7 & 1 & 1 & 0 \\
\bottomrule
\end{tabular}
\end{table}
Παρατηρούμε ότι το \foreignlanguage{english}{Alpha Complex} παράγει πολύ λιγότερα απλόπλοκα (ακμές και τρίγωνα) σε σχέση με το \foreignlanguage{english}{VR}, αλλά διατηρεί τους ίδιους αριθμούς \foreignlanguage{english}{Betti}. Αυτό το καθιστά πολύ πιο αποδοτικό.

\begin{figure}[H]
\centering
\foreignlanguage{english}{\includegraphics[width=0.6\textwidth]{complex_2d_colored.png}}
\caption{Οπτικοποίηση του \foreignlanguage{english}{VR Complex}. Τα σημεία είναι μπλε, οι ακμές κόκκινες και τα τρίγωνα πράσινα.}
\end{figure}

\subsubsection{Α.3—\foreignlanguage{english}{Filtration} \& Ανάλυση Κλίμακας}
\vspace{0.5em}

\vspace{1em}
\textbf{Αποτελέσματα \& Ερμηνεία:}
Η τρύπα του κύκλου γεννιέται στο $r=0.2$ (όταν οι ακμές κλείνουν τον κύκλο) και πεθαίνει στο $r=0.9$ (όταν τα τρίγωνα γεμίζουν τον κύκλο).

\begin{figure}[H]
\centering
\foreignlanguage{english}{\texttt{\includegraphics[width=0.6\textwidth]{filtration_betti.png}}}
\caption{Διάγραμμα αριθμών \foreignlanguage{english}{Betti}. Στο $r=0.2$ εμφανίζεται ο $\beta_1=1$.}
\end{figure}

\begin{figure}[H]
\centering
\foreignlanguage{english}{\includegraphics[width=0.8\textwidth]{complex_steps.png}}
\caption{Στιγμιότυπα της κατασκευής του συμπλέγματος για διάφορες τιμές του $r$.}
\end{figure}

\subsubsection{Α.4—Οπτικοποίηση σε 3 Διαστάσεις (\foreignlanguage{english}{3D})}
\vspace{0.5em}

\vspace{1em}
\textbf{Εξήγηση Υλοποίησης:}
Για τις δομές σε 3 διαστάσεις (όπως ο Τόρος και η Σφαίρα), η οπτικοποίηση του συμπλέγματος \foreignlanguage{english}{Vietoris--Rips} απαιτεί την εμφάνιση των σημείων, των ακμών αλλά και των τριγώνων (\foreignlanguage{english}{2-simplices}) στον τρισδιάστατο χώρο. Αυτό βοηθά στην κατανόηση του πώς τα απλόπλοκα "γεμίζουν" τον χώρο και σχηματίζουν τις τοπολογικές τρύπες.

\begin{figure}[H]
\centering
\foreignlanguage{english}{\includegraphics[width=0.45\textwidth]{complex_3d_persp1.png}}
\foreignlanguage{english}{\includegraphics[width=0.45\textwidth]{complex_3d_persp2.png}}
\caption{Οπτικοποίηση του συμπλέγματος σε 3 διαστάσεις (\foreignlanguage{english}{3D}) για τον Τόρο από δύο διαφορετικές οπτικές γωνίες.}
\end{figure}

\newpage
\subsection{Ερώτημα Β}
\vspace{1em}

\subsubsection{Β.1: \foreignlanguage{english}{Persistent Homology} από Νέφη Σημείων}
\vspace{0.5em}

\vspace{1em}
\textbf{Εξήγηση Υλοποίησης:}
Χρησιμοποιήσαμε τη βιβλιοθήκη \foreignlanguage{english}{\texttt{Ripser}} για τον υπολογισμό των διαγραμμάτων εμμονής. Η \foreignlanguage{english}{\texttt{Ripser}} είναι εξαιρετικά γρήγορη και αποδοτική για αυτόν τον σκοπό.

\vspace{1em}
\textbf{Σημαντικό Κομμάτι Κώδικα:}
\selectlanguage{english}
\begin{lstlisting}
from ripser import ripser
# Calculate Persistent Homology
result = ripser(data, maxdim=2)
diagrams = result['dgms']
\end{lstlisting}
\selectlanguage{greek}
\textit{Σχολιασμός: Η συνάρτηση \foreignlanguage{english}{\texttt{ripser}} υπολογίζει αυτόματα τα διαγράμματα εμμονής μέχρι τη διάσταση που ορίζουμε (\foreignlanguage{english}{\texttt{maxdim}}).}

\vspace{1em}
\textbf{Αποτελέσματα \& Ερμηνεία:}
Τα σημεία που βρίσκονται μακριά από τη διαγώνιο αντιπροσωπεύουν πραγματικά τοπολογικά χαρακτηριστικά που επιμένουν, ενώ τα σημεία κοντά στη διαγώνιο είναι θόρυβος.

\begin{figure}[H]
\centering
\foreignlanguage{english}{\includegraphics[width=0.45\textwidth]{b1_Random_R2.png}}
\foreignlanguage{english}{\includegraphics[width=0.45\textwidth]{b1_Noisy_Circle.png}}
\caption{\foreignlanguage{english}{Persistence Diagrams} για Τυχαία Σημεία και Θορυβώδη Κύκλο.}
\end{figure}

\begin{figure}[H]
\centering
\foreignlanguage{english}{\includegraphics[width=0.45\textwidth]{b1_Torus_T2.png}}
\foreignlanguage{english}{\includegraphics[width=0.45\textwidth]{b1_Sphere_S2.png}}
\caption{\foreignlanguage{english}{Persistence Diagrams} για Τόρο και Σφαίρα.}
\end{figure}

\subsubsection{Β.2: Στιβαρότητα (\foreignlanguage{english}{Robustness}) σε \foreignlanguage{english}{Gaussian} Θόρυβο}
\vspace{0.5em}

\vspace{1em}
\textbf{Αποτελέσματα \& Ερμηνεία:}
Όπως φαίνεται στο διάγραμμα, η απόσταση \foreignlanguage{english}{Bottleneck} και \foreignlanguage{english}{Wasserstein} αυξάνεται καθώς αυξάνεται ο θόρυβος $\sigma$. Αυτό δείχνει την ευαισθησία των τοπολογικών χαρακτηριστικών στον θόρυβο.

\begin{figure}[H]
\centering
\foreignlanguage{english}{\includegraphics[width=0.6\textwidth]{b2_robustness.png}}
\caption{Επίδραση του θορύβου στις αποστάσεις \foreignlanguage{english}{Bottleneck} και \foreignlanguage{english}{Wasserstein}.}
\end{figure}

\subsubsection{Β.3: \foreignlanguage{english}{Persistent Homology} σε Γράφους}
\vspace{0.5em}

\vspace{1em}
\textbf{Αποτελέσματα \& Ερμηνεία:}
Υπολογίστηκαν δύο διαφορετικές διηθήσεις (\foreignlanguage{english}{filtrations}). Τα αποτελέσματα δείχνουν πώς η επιλογή της διήθησης επηρεάζει τα τοπολογικά χαρακτηριστικά που ανιχνεύονται.

\begin{figure}[H]
\centering
\foreignlanguage{english}{\includegraphics[width=0.8\textwidth]{b3_persistence_diagrams_comparison.png}}
\caption{Σύγκριση \foreignlanguage{english}{Persistence Diagrams} για γράφους.}
\end{figure}

\begin{figure}[H]
\centering
\foreignlanguage{english}{\includegraphics[width=0.8\textwidth]{b3_persistence_images.png}}
\caption{\foreignlanguage{english}{Persistence Images} για τους γράφους.}
\end{figure}

\newpage
\subsection{Ερώτημα Γ}
\vspace{1em}

\subsubsection{Γ.1: Υλοποίηση \foreignlanguage{english}{Mapper} εκ του Μηδενός}
\vspace{0.5em}

\vspace{1em}
\textbf{Εξήγηση Υλοποίησης:}
Υλοποιήσαμε τον αλγόριθμο \foreignlanguage{english}{Mapper} από το μηδέν. Ο αλγόριθμος χωρίζει τον χώρο των τιμών του φίλτρου σε επικαλυπτόμενα διαστήματα, ομαδοποιεί τα σημεία που αντιστοιχούν σε κάθε διάστημα (χρησιμοποιώντας \foreignlanguage{english}{DBSCAN}) και συνδέει τους κόμβους που μοιράζονται κοινά σημεία.

\vspace{1em}
\textbf{Σημαντικό Κομμάτι Κώδικα:}
\selectlanguage{english}
\begin{lstlisting}
# Main logic of fit method in CustomMapper
for int_idx, interval_bounds in enumerate(intervals):
    # Find points in interval
    pts_idx = np.where(in_interval)[0]
    # Cluster subset of points
    labels = self.clustering_algo.fit_predict(X_subset)
    # Add nodes and edges to graph
\end{lstlisting}
\selectlanguage{greek}
\textit{Σχολιασμός: Η λογική αυτή διατρέχει όλα τα διαστήματα και εφαρμόζει \en{clustering} για να βρει τους κόμβους του γραφήματος.}

\subsubsection{Γ.2: Εφαρμογή \foreignlanguage{english}{Mapper} \& Ανάλυση Ευαισθησίας}
\vspace{0.5em}

\vspace{1em}
\textbf{Αποτελέσματα \& Ερμηνεία:}
Εφαρμόσαμε τον αλγόριθμο στα σύνολα δεδομένων \foreignlanguage{english}{Figure-8} και \foreignlanguage{english}{Horse-shoe} με 3 διαφορετικές συναρτήσεις φίλτρου: \foreignlanguage{english}{PCA} 1ης συνιστώσας, \foreignlanguage{english}{L2-norm} και \foreignlanguage{english}{Density}.

\begin{figure}[H]
\centering
\foreignlanguage{english}{\includegraphics[width=0.32\textwidth]{g2_figure_8_pca_1st_component_custom.png}}
\foreignlanguage{english}{\includegraphics[width=0.32\textwidth]{g2_figure_8_l2-norm_custom.png}}
\foreignlanguage{english}{\includegraphics[width=0.32\textwidth]{g2_figure_8_density_custom.png}}
\caption{Αποτελέσματα \foreignlanguage{english}{Custom Mapper} για το \foreignlanguage{english}{Figure-8} με φίλτρα \foreignlanguage{english}{PCA}, \foreignlanguage{english}{L2-norm} και \foreignlanguage{english}{Density}.}
\end{figure}

\textit{Σχολιασμός: Η τριάδα των γραφημάτων για το \foreignlanguage{english}{Figure-8} διαφέρει ανάλογα με το φίλτρο:
1. Το \foreignlanguage{english}{PCA} διαχωρίζει τους δύο κύκλους (αριστερά-δεξιά) στον άξονα μέγιστης διασποράς.
2. Το \foreignlanguage{english}{L2-norm} δείχνει την απόσταση από το κέντρο, δημιουργώντας μια δομή που "ανοίγει" από το κέντρο προς τα άκρα.
3. Το \foreignlanguage{english}{Density} αναδεικνύει το σημείο τομής των δύο κύκλους όπου η πυκνότητα των σημείων είναι μέγιστη.}

\begin{figure}[H]
\centering
\foreignlanguage{english}{\includegraphics[width=0.32\textwidth]{g2_horse_shoe_pca_1st_component_custom.png}}
\foreignlanguage{english}{\includegraphics[width=0.32\textwidth]{g2_horse_shoe_l2-norm_custom.png}}
\foreignlanguage{english}{\includegraphics[width=0.32\textwidth]{g2_horse_shoe_density_custom.png}}
\caption{Αποτελέσματα \foreignlanguage{english}{Custom Mapper} για το \foreignlanguage{english}{Horse-shoe} με φίλτρα \foreignlanguage{english}{PCA}, \foreignlanguage{english}{L2-norm} και \foreignlanguage{english}{Density}.}
\end{figure}

\textit{Σχολιασμός: Στο \foreignlanguage{english}{Horse-shoe}:
1. Το \foreignlanguage{english}{PCA} "ξεδιπλώνει" την καμπύλη σε μια συνεχή γραμμική δομή.
2. Το \foreignlanguage{english}{L2-norm} δίνει παρόμοιες τιμές σε όλα τα σημεία (αφού ισαπέχουν από το κέντρο), συμπυκνώνοντας το γράφημα.
3. Το \foreignlanguage{english}{Density} δείχνει ότι η πυκνότητα είναι ομοιόμορφη κατά μήκος του πετάλου, διατηρώντας τη δομή ως μία συνεκτική συνιστώσα.}

\begin{figure}[H]
\centering
\foreignlanguage{english}{\includegraphics[width=0.6\textwidth]{g2_sensitivity_heatmap.png}}
\caption{Ανάλυση ευαισθησίας: Αριθμός συνεκτικών συνιστωσών συναρτήσει των παραμέτρων.}
\end{figure}

\textit{Σχολιασμός: Η ανάλυση ευαισθησίας δείχνει ότι ο αριθμός των συνεκτικών συνιστωσών εξαρτάται έντονα από τον αριθμό των διαστημάτων και το ποσοστό επικάλυψης. Μεγαλύτερη επικάλυψη οδηγεί σε πιο συνεκτικά γραφήματα.}

\newpage
\subsection{Ερώτημα Δ --- Εφαρμογή σε Πραγματικά Δεδομένα}
\vspace{1em}

Σε αυτή την ενότητα εφαρμόζουμε τις τεχνικές της Τοπολογικής Ανάλυσης Δεδομένων (\foreignlanguage{english}{TDA}) σε δύο πραγματικά σύνολα δεδομένων: το \foreignlanguage{english}{dataset} καρκίνου του μαστού (\foreignlanguage{english}{Breast Cancer Wisconsin}) ως Κατηγορία Ι και τις χρονοσειρές των τιμών κλεισίματος του δείκτη \foreignlanguage{english}{S\&P 500} των τελευταίων δύο ετών ως Κατηγορία ΙΙ.

\subsubsection{Δ.1 --- Ανάλυση Πρώτου \en{Dataset} (\en{Breast Cancer})}
\vspace{0.5em}

\textbf{1. Προεπεξεργασία και Μείωση Διάστασης:}
Το σύνολο δεδομένων \foreignlanguage{english}{Breast Cancer} περιλαμβάνει 569 δείγματα με 30 αριθμητικά χαρακτηριστικά που περιγράφουν πυρήνες κυττάρων από βιοψίες μαστού. Τα δεδομένα κατηγοριοποιούνται σε δύο κλάσεις: \foreignlanguage{english}{Malignant} (κακοήθης, 212 δείγματα) και \foreignlanguage{english}{Benign} (καλοήθης, 357 δείγματα). 
Πραγματοποιήθηκε τυποποίηση (\foreignlanguage{english}{Standardization}) των δεδομένων με χρήση του \foreignlanguage{english}{\texttt{StandardScaler}} ώστε όλα τα χαρακτηριστικά να έχουν μέση τιμή 0 και διασπορά 1. Στη συνέχεια, εφαρμόστηκε η μέθοδος \foreignlanguage{english}{PCA} κρατώντας τις 5 κύριες συνιστώσες για τη διευκόλυνση του υπολογισμού της επίμονης ομολογίας.

\textbf{2. Επίμονη Ομολογία ($H_0, H_1$):}
Υπολογίστηκε η επίμονη ομολογία για τις διαστάσεις 0 και 1 ξεχωριστά για την κλάση των κακοήθων και των καλοήθων δειγμάτων χρησιμοποιώντας τη βιβλιοθήκη \foreignlanguage{english}{\texttt{ripser}}. 
Τα αποτελέσματα οπτικοποιούνται στα διαγράμματα εμμονής (\foreignlanguage{english}{Persistence Diagrams}) και \foreignlanguage{english}{barcodes}.

\begin{figure}[H]
\centering
\foreignlanguage{english}{\includegraphics[width=0.8\textwidth]{d1_persistence_diagrams.png}}
\caption{\foreignlanguage{english}{Persistence Diagrams} για την κλάση \foreignlanguage{english}{Malignant} (αριστερά) και \foreignlanguage{english}{Benign} (δεξιά).}
\end{figure}

Παρατηρούμε ότι οι δύο κλάσεις παρουσιάζουν διαφορές στην τοπολογική τους δομή. Η κλάση των καλοήθων δειγμάτων εμφανίζει πιο διάσπαρτα και μακρόβια χαρακτηριστικά $H_1$ (τρύπες) σε σχέση με την κλάση των κακοήθων, πράγμα που υποδηλώνει μια πιο σύνθετη και «ανοιχτή» γεωμετρική διάταξη στο χώρο των κύριων συνιστωσών.

\textbf{3. Εφαρμογή \foreignlanguage{english}{Mapper}:}
Κατασκευάστηκε το γράφημα \foreignlanguage{english}{Mapper} χρησιμοποιώντας τη βιβλιοθήκη \foreignlanguage{english}{\texttt{kmapper}}. Ως συνάρτηση φίλτρου χρησιμοποιήθηκε η προβολή \foreignlanguage{english}{PCA} σε 2 διαστάσεις. Ο χώρος καλύφθηκε από 10 επικαλυπτόμενα διαστήματα με ποσοστό επικάλυψης 30\% ανά διάσταση, και εφαρμόστηκε ο αλγόριθμος \foreignlanguage{english}{DBSCAN} (με \foreignlanguage{english}{\texttt{eps=2.5}} και \foreignlanguage{english}{\texttt{min\_samples=3}}) για το \foreignlanguage{english}{clustering} των σημείων. Το διαδραστικό γράφημα \foreignlanguage{english}{Mapper} αποθηκεύτηκε στο αρχείο \foreignlanguage{english}{\texttt{d1\_mapper\_graph.html}}. Το γράφημα αναδεικνύει τον διαχωρισμό των κλάσεων σε διαφορετικές περιοχές (\foreignlanguage{english}{branches}) του γραφήματος, με μια περιοχή μετάβασης όπου συνυπάρχουν καλοήθη και κακοήθη δείγματα.

\textbf{4. \foreignlanguage{english}{TDA} ως Χαρακτηριστικά (\foreignlanguage{english}{Features}) για Μηχανική Μάθηση:}
Για να χρησιμοποιήσουμε την τοπολογική πληροφορία σε ταξινόμηση, μετατρέψαμε το διάνυσμα χαρακτηριστικών κάθε ασθενούς (30 τιμές) σε χρονοσειρά / \foreignlanguage{english}{point cloud} μέσω της μεθόδου \foreignlanguage{english}{Takens Embedding} (με διάσταση 3 και καθυστέρηση 1).
Στη συνέχεια, για κάθε ασθενή υπολογίστηκε η επίμονη ομολογία $H_0$ και τα διαγράμματα εμμονής μετατράπηκαν σε εικόνες εμμονής (\foreignlanguage{english}{Persistence Images}) μέσω του \foreignlanguage{english}{\texttt{PersistenceImager}} (με μέγεθος \foreignlanguage{english}{pixel} 0.1). Οι εικόνες αυτές ισοπεδώθηκαν (\foreignlanguage{english}{flatten}) σε μονοδιάστατα διανύσματα χαρακτηριστικών.
Εκπαιδεύτηκε ένας ταξινομητής \foreignlanguage{english}{Random Forest} με τα παρακάτω αποτελέσματα στο \foreignlanguage{english}{test set} (30\%):
\begin{itemize}
    \item \foreignlanguage{english}{\textbf{Accuracy}}: 0.6901
    \item \foreignlanguage{english}{\textbf{AUC Score}}: 0.7152
\end{itemize}
Αν και η ακρίβεια αυτή είναι χαμηλότερη από τη χρήση των αρχικών δεδομένων άμεσα (καθώς το \foreignlanguage{english}{Takens embedding} σε 30 \foreignlanguage{english}{features} εισάγει θόρυβο), επιβεβαιώνει ότι η καθαρά τοπολογική δομή των χαρακτηριστικών φέρει σημαντική διακριτική ικανότητα για τη διάγνωση του καρκίνου.

\subsubsection{Δ.2 --- Ανάλυση Δεύτερου \en{Dataset} (\en{S\&P 500})}
\vspace{0.5em}

\textbf{1. Δεδομένα και \foreignlanguage{english}{Giotto-TDA Pipeline}:}
Φορτώθηκαν οι καθημερινές τιμές κλεισίματος του δείκτη \foreignlanguage{english}{S\&P 500} (\foreignlanguage{english}{\texttt{\^{}GSPC}}) για τα τελευταία 2 έτη μέσω της \foreignlanguage{english}{\texttt{yfinance}}. Οι τιμές τυποποιήθηκαν και εφαρμόστηκε η μέθοδος \foreignlanguage{english}{Takens Embedding} με διάσταση $d=3$ και καθυστέρηση $\tau=1$, μετατρέποντας τη μονοδιάστατη χρονοσειρά σε ένα τρισδιάστατο \foreignlanguage{english}{point cloud} που αναπαριστά τη δυναμική της αγοράς.
Κατασκευάστηκε ένα \foreignlanguage{english}{scikit-learn} συμβατό \foreignlanguage{english}{pipeline} μέσω της βιβλιοθήκης \foreignlanguage{english}{\texttt{giotto-tda}} που περιλαμβάνει:
\begin{enumerate}
    \item \foreignlanguage{english}{\texttt{TakensEmbedding}} για την αναπαράσταση της χρονοσειράς σε \foreignlanguage{english}{3D}.
    \item \foreignlanguage{english}{\texttt{VietorisRipsPersistence}} για τον υπολογισμό των διαγραμμάτων εμμονής για τις διαστάσεις 0, 1 και 2.
    \item \foreignlanguage{english}{\texttt{PersistenceEntropy}} για τον υπολογισμό της επίμονης εντροπίας ως μέτρο πολυπλοκότητας της τοπολογίας της χρονοσειράς.
\end{enumerate}

Η κανονικοποιημένη επίμονη εντροπία (\foreignlanguage{english}{Normalized Persistence Entropy}) για το σύνολο των δεδομένων υπολογίστηκε σε:
\begin{itemize}
    \item $H_0$: 1.6775
    \item $H_1$: 8.6782
    \item $H_2$: -0.6358
\end{itemize}

\textbf{2. Επαλήθευση $H_2$ Ομολογίας με \foreignlanguage{english}{GUDHI}:}
Επειδή η διάσταση του \foreignlanguage{english}{point cloud} είναι $d \ge 3$, υπολογίστηκε η ομολογία διάστασης 2 ($H_2$). Χρησιμοποιήθηκε το \foreignlanguage{english}{\texttt{AlphaComplex}} της βιβλιοθήκης \foreignlanguage{english}{GUDHI} για την επαλήθευση. Οι αριθμοί \foreignlanguage{english}{Betti} που προέκυψαν ήταν $\beta = [1, 0, 0]$, γεγονός που δείχνει ότι το σύμπλεγμα είναι τοπολογικά ισοδύναμο με ένα μόνο σημείο (συνεκτικό σύμπλεγμα χωρίς 1-κύκλους ή 2-κοιλότητες στην τελική κλίμακα), επιβεβαιώνοντας ότι το \foreignlanguage{english}{Takens Embedding} της χρηματοοικονομικής χρονοσειράς δεν σχηματίζει κλειστές σφαιρικές κοιλότητες ($H_2 = 0$).

\textbf{3. Οπτικοποίηση και Μείωση Διάστασης:}
Η \foreignlanguage{english}{3D} τροχιά του \foreignlanguage{english}{Takens Embedding} οπτικοποιήθηκε σε διαδραστικό γράφημα \foreignlanguage{english}{Plotly} (\foreignlanguage{english}{\texttt{d2\_sp500\_3d\_embedding.html}}). Επιπλέον, το \foreignlanguage{english}{3D point cloud} οπτικοποιήθηκε και σε στατικό τρισδιάστατο γράφημα (\foreignlanguage{english}{embedding}).

\begin{figure}[H]
\centering
\foreignlanguage{english}{\includegraphics[width=0.6\textwidth]{d2_sp500_3d_embedding.png}}
\caption{Τρισδιάστατη τροχιά του \foreignlanguage{english}{Takens Embedding} ($d=3, \tau=1$) για τη χρονοσειρά \foreignlanguage{english}{S\&P 500}.}
\end{figure}

Για τη σύγκριση, το \foreignlanguage{english}{3D point cloud} προβλήθηκε σε 2 διαστάσεις με 3 μεθόδους: \foreignlanguage{english}{PCA}, \foreignlanguage{english}{t-SNE} και \foreignlanguage{english}{UMAP}.

\begin{figure}[H]
\centering
\foreignlanguage{english}{\includegraphics[width=0.9\textwidth]{d2_sp500_dim_reduction.png}}
\caption{Σύγκριση μείωσης διάστασης του \foreignlanguage{english}{3D Takens Embedding} του \foreignlanguage{english}{S\&P 500} με \foreignlanguage{english}{PCA}, \foreignlanguage{english}{t-SNE} και \foreignlanguage{english}{UMAP}.}
\end{figure}

\begin{itemize}
    \item Η \foreignlanguage{english}{PCA} αποκαλύπτει τη γραμμική κατεύθυνση μέγιστης διασποράς της αγοράς.
    \item Το \foreignlanguage{english}{t-SNE} δημιουργεί ξεχωριστές συστάδες, τονίζοντας τοπικές ομοιότητες αλλά καταστρέφοντας την παγκόσμια γεωμετρία της χρονοσειράς.
    \item Το \foreignlanguage{english}{UMAP} διατηρεί καλύτερα τη συνέχεια της τροχιάς της χρονοσειράς, αναδεικνύοντας τη δομή «νήματος» (\foreignlanguage{english}{filament structure}) των τιμών.
\end{itemize}

\textbf{4. \foreignlanguage{english}{Bootstrap Resampling} \& Διαστήματα Εμπιστοσύνης (95\%):}
Για να ελέγξουμε τη σταθερότητα της επίμονης εντροπίας, εφαρμόστηκε \foreignlanguage{english}{bootstrap resampling} ($n=100$) με τυχαία επιλογή του 80\% των σημείων με επανατοποθέτηση. Τα 95\% διαστήματα εμπιστοσύνης για την επίμονη εντροπία είναι:
\begin{itemize}
    \item \textbf{$H_0$}: \foreignlanguage{english}{Mean} = 1.7014, 95\% \foreignlanguage{english}{CI} = $[1.6596, 1.7349]$
    \item \textbf{$H_1$}: \foreignlanguage{english}{Mean} = -12.2571, 95\% \foreignlanguage{english}{CI} = $[-28.6147, -4.8745]$
    \item \textbf{$H_2$}: \foreignlanguage{english}{Mean} = -0.2339, 95\% \foreignlanguage{english}{CI} = $[-1.0000, -0.0000]$
\end{itemize}
Τα στενά διαστήματα εμπιστοσύνης για το $H_0$ δείχνουν τη σταθερότητα της βασικής συνεκτικότητας των δεδομένων, ενώ οι αρνητικές τιμές και η διακύμανση στο $H_1$ και $H_2$ υποδηλώνουν την ευαισθησία και τον θόρυβο των ανώτερων τοπολογικών χαρακτηριστικών όταν υποδειγματίζεται το \foreignlanguage{english}{point cloud} της χρονοσειράς.

\subsubsection{Δ.3 --- Βαθιά Σύγκριση \& Συζήτηση: \en{TDA} \en{vs} Κλασικές Μέθοδοι}
\vspace{0.5em}

Η εφαρμογή των \foreignlanguage{english}{TDA} μεθόδων στα δύο \foreignlanguage{english}{datasets} ανέδειξε τα εξής πλεονεκτήματα έναντι των \foreignlanguage{english}{PCA}, \foreignlanguage{english}{t-SNE} και \foreignlanguage{english}{UMAP}:
\begin{enumerate}
    \item \textbf{Πολλαπλές Κλίμακες:} Οι κλασικές μέθοδοι μείωσης διάστασης απαιτούν την επιλογή μιας συγκεκριμένης κλίμακας (π.χ. \foreignlanguage{english}{perplexity} στο \foreignlanguage{english}{t-SNE} ή \foreignlanguage{english}{\texttt{n\_neighbors}} στο \foreignlanguage{english}{UMAP}). Αντίθετα, η επίμονη ομολογία εξετάζει όλες τις κλίμακες ταυτόχρονα (μέσω του \en{filtration} $r$), διαχωρίζοντας με μαθηματική αυστηρότητα το σήμα από τον θόρυβο.
    \item \foreignlanguage{english}{\textbf{Feature Extraction}}: Η \foreignlanguage{english}{TDA} επιτρέπει τη μετατροπή γεωμετρικών ιδιοτήτων σε ποσοτικά χαρακτηριστικά (όπως η \foreignlanguage{english}{Persistent Entropy} και οι \foreignlanguage{english}{Persistence Images}) που εισάγονται άμεσα σε μοντέλα μηχανικής μάθησης. Οι κλασικές μέθοδοι (όπως το \foreignlanguage{english}{t-SNE}) δεν ορίζουν εύκολα μια αντίστροφη απεικόνιση ή μια γενικεύσιμη μεταμόρφωση για νέα δεδομένα (\foreignlanguage{english}{out-of-sample extension}).
    \item \textbf{Διατήρηση Σχήματος:} Ο \foreignlanguage{english}{Mapper} και η επίμονη εντροπία αποτυπώνουν το «σχήμα» των δεδομένων (π.χ. βρόχους, διακλαδώσεις, κοιλότητες) το οποίο οι κλασικές προβολές (κυρίως η \foreignlanguage{english}{PCA}) συμπιέζουν και καταστρέφουν. Στην ανάλυση χρονοσειρών (\foreignlanguage{english}{S\&P 500}), η \foreignlanguage{english}{TDA} σε συνδυασμό με το \foreignlanguage{english}{Takens Embedding} προσφέρει μια πλήρη εικόνα του χώρου φάσεων (\foreignlanguage{english}{phase space reconstruction}), επιτρέποντας την ανίχνευση τοπολογικών αλλαγών που σχετίζονται με περιόδους αστάθειας στην αγορά.
\end{enumerate}

\newpage
\subsection{Ερώτημα Ε --- \en{Zigzag Persistent Homology}}
\vspace{1em}

\subsubsection{Ε.1 --- \en{Zigzag Persistent Homology} σε Χρονοσειρές}
\vspace{0.5em}

\textbf{1. Μεθοδολογία και Κατασκευή:}
Για την καταγραφή της δυναμικής εξέλιξης των τοπολογικών χαρακτηριστικών της χρονοσειράς του δείκτη \en{S\&P 500}, εφαρμόσαμε τη μέθοδο της \en{Zigzag Persistent Homology}. Η μονοδιάστατη χρονοσειρά μετατράπηκε σε ακολουθία \en{point clouds} μέσω \en{Sliding Window} με μέγεθος παραθύρου $W = 50$ ημέρες και βήμα $S = 25$ ημέρες, χρησιμοποιώντας \en{Takens Embedding} με διάσταση $d = 3$ και καθυστέρηση $\tau = 1$. 

Η κατασκευή της \en{zigzag filtration} έγινε με την εναλλαγή \en{forward} και \en{backward inclusions}:
\[ K_0 \hookrightarrow K_0 \cup K_1 \hookleftarrow K_1 \hookrightarrow K_1 \cup K_2 \hookleftarrow K_2 \hookrightarrow \dots \]
όπου $K_i$ είναι το \en{Rips complex} του $i$-οστού παραθύρου για ακτίνα \en{Rips} $\epsilon = 2 \times 0.1524 = 0.3048$ (η 15η εκατοστημόρια θέση των αποστάσεων). Ο υπολογισμός των αριθμών \en{Betti} έγινε μέσω απαλοιφής \en{Gauss} σε \en{GF(2)} πάνω στους αντίστοιχους πίνακες συνοριακών απεικονίσεων ($\partial_1, \partial_2$).

\begin{figure}[H]
\centering
\foreignlanguage{english}{\includegraphics[width=0.85\textwidth]{e1_sliding_windows.png}}
\caption{\en{Sliding windows} στη χρονοσειρά \en{S\&P 500} που χρησιμοποιήθηκαν για την κατασκευή των \en{point clouds}.}
\label{fig:e1_sliding_windows}
\end{figure}

\textbf{2. Αποτελέσματα και Σχολιασμός:}
\begin{itemize}
    \item \textbf{Συνεκτικότητα ($H_0$):} Εντοπίστηκαν $13$ \en{intervals} στην $H_0$. Το κυρίαρχο \en{interval} έχει διάρκεια $19.0$ \en{steps} (από $t=0$ έως $t=19$), που αντιστοιχεί στη βασική συνεκτική δομή που επιβιώνει σε όλα τα παράθυρα. Παρατηρείται ένα σημαντικό \en{spike} στο $t = 16$ (παράθυρο $K_8$), όπου ο αριθμός των συνεκτικών συνιστωσών αυξάνεται σε $\beta_0 = 9$, υποδηλώνοντας τοπολογική διάσπαση του \en{point cloud} κατά τη διάρκεια περιόδου υψηλής μεταβλητότητας στην αγορά.
    \item \textbf{Τοπολογικοί Κύκλοι ($H_1$):} Εντοπίστηκε $1$ \en{interval} στην $H_1$ με διάρκεια $1.0$ \en{step}. Η κλασική \en{PH} ανά παράθυρο δίνει μέσο $\beta_1 = 0.00 \pm 0.00$. Αυτό δείχνει ότι τα μεμονωμένα παράθυρα δεν περιέχουν μόνιμους κύκλους, αλλά η \en{zigzag PH} ανιχνεύει έναν μεταβατικό (\en{transient}) κύκλο που δημιουργείται κατά την ένωση $K_i \cup K_{i+1}$.
\end{itemize}

\begin{figure}[H]
\centering
\foreignlanguage{english}{\includegraphics[width=0.9\textwidth]{e1_summary_comparison.png}}
\caption{Συγκεντρωτικό διάγραμμα σύγκρισης: \en{Zigzag Persistence Diagram} (πάνω αριστερά), \en{Zigzag Barcode} (πάνω δεξιά) και εξέλιξη των \en{Betti} αριθμών $\beta_0$ (κάτω αριστερά) και $\beta_1$ (κάτω δεξιά) σε σύγκριση με την κλασική \en{PH}.}
\label{fig:e1_summary_comparison}
\end{figure}

\textbf{3. Σύγκριση με Κλασική \en{Persistent Homology}:}
Η κλασική \en{PH} αναλύει κάθε παράθυρο ανεξάρτητα και στατικά, χάνοντας τη χρονική συσχέτιση των χαρακτηριστικών. Αντίθετα, η \en{zigzag PH} συνδέει τα τοπολογικά χαρακτηριστικά μεταξύ διαδοχικών παραθύρων μέσω των \en{union complexes}. 

Όπως επιβεβαιώνεται στα διαγράμματα του Σχήματος \ref{fig:e1_summary_comparison}, στα άρτια βήματα $t=2i$ οι Betti αριθμοί της \en{zigzag PH} ταυτίζονται απόλυτα με τους \en{Betti} αριθμούς της κλασικής \en{PH} στο παράθυρο $i$, ενώ τα ενδιάμεσα περιττά βήματα αποτυπώνουν με ακρίβεια τη δυναμική μετάβαση (ένωση) μεταξύ τους.

\newpage
\section{\texorpdfstring{\underline{Συμπεράσματα}}{Συμπεράσματα}}
\vspace{1em}
Η \foreignlanguage{english}{TDA} αποκάλυψε ότι τα δεδομένα μας διαθέτουν πλούσια τοπολογική δομή που δεν είναι πάντα ορατή με παραδοσιακές μεθόδους στατιστικής ή μηχανικής μάθησης. Συγκεκριμένα:

\begin{enumerate}
    \item Οι αριθμοί \foreignlanguage{english}{Betti} και τα διαγράμματα εμμονής προσδιορίζουν με ακρίβεια τον αριθμό των οπών και των κοιλοτήτων στα σχήματα (κύκλος, τόρος, σφαίρα).
    \item Η ανάλυση ευαισθησίας στον θόρυβο έδειξε ότι τα τοπολογικά χαρακτηριστικά είναι στιβαρά μέχρι ένα σημείο, αλλά ο έντονος θόρυβος μπορεί να αλλοιώσει τις αποστάσεις \foreignlanguage{english}{Bottleneck}.
    \item Ο αλγόριθμος \foreignlanguage{english}{Mapper} αποδείχθηκε εξαιρετικό εργαλείο για την οπτικοποίηση σύνθετων δομών όπως το \foreignlanguage{english}{Figure-8} και το \foreignlanguage{english}{Horse-shoe}, επιτρέποντας την κατανόηση της γεωμετρίας τους μέσω απλών γραφημάτων.
    \item Η μέθοδος της \foreignlanguage{english}{Zigzag Persistent Homology} επέτρεψε την επιτυχή καταγραφή της δυναμικής εξέλιξης των τοπολογικών δομών σε χρηματοοικονομικές χρονοσειρές, αναδεικνύοντας μεταβατικά τοπολογικά χαρακτηριστικά που συνδέονται με την αστάθεια της αγοράς.
\end{enumerate}

\newpage
\section{\texorpdfstring{\underline{Χρήση \en{AI}}}{Χρήση AI}}
\vspace{1em}
Η χρήση Τεχνητής Νοημοσύνης (συγκεκριμένα Μεγάλων Γλωσσικών Μοντέλων - \foreignlanguage{english}{LLMs}) έγινε στα εξής σημεία της εργασίας:

\begin{itemize}
    \item \textbf{Ανάπτυξη Κώδικα:} Υποστήριξη στη συγγραφή και βελτιστοποίηση κώδικα \foreignlanguage{english}{Python} σε όλα τα ερωτήματα της εργασίας (συμπεριλαμβανομένης της υλοποίησης του συμπλέγματος \foreignlanguage{english}{Vietoris--Rips} εκ του μηδενός και του αλγορίθμου \foreignlanguage{english}{Custom Mapper}).
    \item \textbf{Επίλυση Σφαλμάτων (\foreignlanguage{english}{Debugging}):} Βοήθεια στην επίλυση σφαλμάτων συμβατότητας των βιβλιοθηκών \foreignlanguage{english}{GUDHI} και \foreignlanguage{english}{Ripser}, καθώς και στη διαχείριση της κωδικοποίησης (\foreignlanguage{english}{UTF-8}) στο περιβάλλον των \foreignlanguage{english}{Windows}.
    \item \textbf{Συγγραφή Αναφοράς σε \foreignlanguage{english}{\LaTeX}:} Υποστήριξη στη δομή του εγγράφου, στη διαχείριση των πακέτων γλωσσών (\foreignlanguage{english}{\texttt{babel}}) για τη σωστή εναλλαγή ελληνικών και αγγλικών, και στη μορφοποίηση των πινάκων και των σχημάτων.
    \item \textbf{Ερμηνεία Αποτελεσμάτων:} Συζήτηση και ανάλυση των διαγραμμάτων εμμονής και των γραφημάτων του \foreignlanguage{english}{Mapper} για την καλύτερη κατανόηση των τοπολογικών δομών.
\end{itemize}

\textit{Σημείωση: Η τελική αξιολόγηση, ο έλεγχος της ορθότητας του κώδικα και η συγγραφή των τελικών συμπερασμάτων έγιναν από την ομάδα εργασίας.}

\newpage
\section{\texorpdfstring{\underline{Βιβλιογραφία}}{Βιβλιογραφία}}
\vspace{1em}
\begin{thebibliography}{9}
\bibitem{andrikopoulos} Ανδρικόπουλος, Α. \emph{Τοπολογική Ανάλυση Δεδομένων}. Πανεπιστήμιο Πατρών, Τμήμα Μηχανικών Η/Υ και Πληροφορικής.
\bibitem{gudhi} \foreignlanguage{english}{GUDHI Editorial Board}. \emph{\foreignlanguage{english}{GUDHI User and Reference Manual}}.
\bibitem{ripser} \foreignlanguage{english}{Bauer, U.} (2021). \emph{\foreignlanguage{english}{Ripser: efficient computation of Vietoris-Rips persistence barcodes}}.
\bibitem{carlsson} Carlsson, G., \& de Silva, V. (2010). \foreignlanguage{english}{Zigzag Persistence}. \emph{\foreignlanguage{english}{Foundations of Computational Mathematics}}, 10(4), 367-405.
\end{thebibliography}

\end{document}
